% GENERATED FILE. Edit canonical lessons, then run npm run generate:books.
% canonical-source-hash: fnv1a64:8d93284bfa1e734a
% canonical-lessons: JA-C377-chien, JA-C377-unkyuu, JA-C377-jikokuhyou, JA-C377-houmen, JA-C377-hanareru, JA-R377-review-372-377

\chapter{Delays and Cancellations}
\label{ch:ja-delays-and-cancellations}
\hlchaptermodality{377}

\begin{chapteropening}
\textbf{By the end of this chapter:} \emph{I can understand why a train is late or not running.}

Complete the last lesson of chapter 377: I can understand why a train is late or not running.
\end{chapteropening}

\section[\texorpdfstring{\ja{ちえん} (chien)}{ちえん (chien)}]{\ja{ちえん} (chien) — a delay}

\label{lesson:JA-C377-chien}



\begin{quote}
Before the new one: say the Japanese for a route, then the Japanese for to skip.
\end{quote}

\subsection*{You'll want to know: \ja{ちえん}}
\textbf{\ja{ちえん}} — \emph{chien} — \textquotedblleft{}a delay\textquotedblright{}.

\begin{quote}
\textbf{\ja{たいふうで} \ja{でんしゃに} \ja{ちえんが} \ja{ありました。}} — \emph{taifū de densha ni chien ga arimashita} — \textquotedblleft{}The typhoon delayed the trains.\textquotedblright{}
\end{quote}

\begin{cousinweb}
\textbf{\ja{ちえん}} (\emph{chien}) is Sino-Japanese (kango): two Chinese characters read with their on'yomi, readings borrowed from Chinese. They are \emph{chi}, \textquotedblleft{}late\textquotedblright{} (Middle Chinese \emph{drij}) and \emph{en}, \textquotedblleft{}prolong\textquotedblright{} (\emph{yen}) (KANJIDIC2; Middle Chinese in Baxter's transcription).
\end{cousinweb}

\begin{grammarlens}[title={What you've built}]
One more word for getting around.
\end{grammarlens}

\subsection*{Your turn}
\begin{itemize}
\raggedright
  \item \emph{Say it:} \emph{chien}
  \item \emph{Say it:} \emph{chien}, once more
  \item \emph{Say it:} \emph{taifū de densha ni chien ga arimashita} — The typhoon delayed the trains.
  \item \emph{Recall:} say the Japanese for a route, then the Japanese for to skip, then say \emph{chien} again
\end{itemize}

\begin{tcolorbox}[breakable,colback=teal!4,colframe=teal!35!black,title={Before you move on}]
What does \ja{ちえん} mean? (\textquotedblleft{}A delay\textquotedblright{}.) Say it once more.
\end{tcolorbox}
\section[\texorpdfstring{\ja{うんきゅう} (unkyū)}{うんきゅう (unkyū)}]{\ja{うんきゅう} (unkyū) — a suspension of service}

\label{lesson:JA-C377-unkyuu}



\begin{quote}
Before the new one: say the Japanese for to skip, then the Japanese for a delay.
\end{quote}

\subsection*{You'll want to know: \ja{うんきゅう}}
\textbf{\ja{うんきゅう}} — \emph{unkyū} — \textquotedblleft{}a suspension of service\textquotedblright{}.

\begin{quote}
\textbf{\ja{ゆきで} \ja{しんかんせんが} \ja{うんきゅうに} \ja{なりました。}} — \emph{yuki de shinkansen ga unkyū ni narimashita} — \textquotedblleft{}Because of the snow, the bullet train stopped running.\textquotedblright{}
\end{quote}

\begin{cousinweb}
\textbf{\ja{うんきゅう}} (\emph{unkyū}) is Sino-Japanese (kango): two Chinese characters read with their on'yomi, readings borrowed from Chinese. They are \emph{un}, \textquotedblleft{}transport\textquotedblright{} (Middle Chinese \emph{hjunH}) and \emph{kyū}, \textquotedblleft{}rest\textquotedblright{} (\emph{xjuw}) (KANJIDIC2; Middle Chinese in Baxter's transcription).
\end{cousinweb}

\begin{grammarlens}[title={What you've built}]
One more word for getting around.
\end{grammarlens}

\subsection*{Your turn}
\begin{itemize}
\raggedright
  \item \emph{Say it:} \emph{unkyū}
  \item \emph{Say it:} \emph{unkyū}, once more
  \item \emph{Say it:} \emph{yuki de shinkansen ga unkyū ni narimashita} — Because of the snow, the bullet train stopped running.
  \item \emph{Recall:} say the Japanese for to skip, then the Japanese for a delay, then say \emph{unkyū} again
\end{itemize}

\begin{tcolorbox}[breakable,colback=teal!4,colframe=teal!35!black,title={Before you move on}]
What does \ja{うんきゅう} mean? (\textquotedblleft{}A suspension of service\textquotedblright{}.) Say it once more.
\end{tcolorbox}
\section[\texorpdfstring{\ja{じこくひょう} (jikokuhyō)}{じこくひょう (jikokuhyō)}]{\ja{じこくひょう} (jikokuhyō) — a timetable}

\label{lesson:JA-C377-jikokuhyou}



\begin{quote}
Before the new one: say the Japanese for a delay, then the Japanese for a suspension of service.
\end{quote}

\subsection*{You'll want to know: \ja{じこくひょう}}
\textbf{\ja{じこくひょう}} — \emph{jikokuhyō} — \textquotedblleft{}a timetable\textquotedblright{}.

\begin{quote}
\textbf{\ja{じこくひょうで} \ja{しはつの} \ja{じかんを} \ja{しらべました。}} — \emph{jikokuhyō de shihatsu no jikan o shirabemashita} — \textquotedblleft{}I checked the time of the first train in the timetable.\textquotedblright{}
\end{quote}

\begin{cousinweb}
\textbf{\ja{じこくひょう}} (\emph{jikokuhyō}) is Sino-Japanese (kango): three Chinese characters read with their on'yomi, readings borrowed from Chinese. They are \emph{ji}, \textquotedblleft{}time\textquotedblright{} (Middle Chinese \emph{dzyi}), \emph{koku}, \textquotedblleft{}engrave\textquotedblright{} (\emph{khok}) and \emph{hyō}, \textquotedblleft{}table\textquotedblright{} (\emph{pjewX}) (KANJIDIC2; Middle Chinese in Baxter's transcription).
\end{cousinweb}

\begin{grammarlens}[title={What you've built}]
One more word for getting around.
\end{grammarlens}

\subsection*{Your turn}
\begin{itemize}
\raggedright
  \item \emph{Say it:} \emph{jikokuhyō}
  \item \emph{Say it:} \emph{jikokuhyō}, once more
  \item \emph{Say it:} \emph{jikokuhyō de shihatsu no jikan o shirabemashita} — I checked the time of the first train in the timetable.
  \item \emph{Recall:} say the Japanese for a delay, then the Japanese for a suspension of service, then say \emph{jikokuhyō} again
\end{itemize}

\begin{tcolorbox}[breakable,colback=teal!4,colframe=teal!35!black,title={Before you move on}]
What does \ja{じこくひょう} mean? (\textquotedblleft{}A timetable\textquotedblright{}.) Say it once more.
\end{tcolorbox}
\section[\texorpdfstring{\ja{ほうめん} (hōmen)}{ほうめん (hōmen)}]{\ja{ほうめん} (hōmen) — a direction (of travel); an area}

\label{lesson:JA-C377-houmen}



\begin{quote}
Before the new one: say the Japanese for a suspension of service, then the Japanese for a timetable.
\end{quote}

\subsection*{You'll want to know: \ja{ほうめん}}
\textbf{\ja{ほうめん}} — \emph{hōmen} — \textquotedblleft{}a direction (of travel); an area\textquotedblright{}.

\begin{quote}
\textbf{\ja{みなと} \ja{ほうめんの} \ja{でんしゃに} \ja{のりました。}} — \emph{minato hōmen no densha ni norimashita} — \textquotedblleft{}I took a train going toward the harbour.\textquotedblright{}
\end{quote}

\begin{cousinweb}
\textbf{\ja{ほうめん}} (\emph{hōmen}) is Sino-Japanese (kango): two Chinese characters read with their on'yomi, readings borrowed from Chinese. They are \emph{hō}, \textquotedblleft{}direction\textquotedblright{} (Middle Chinese \emph{pjang}) and \emph{men}, \textquotedblleft{}face\textquotedblright{} (\emph{mjienH}) (KANJIDIC2; Middle Chinese in Baxter's transcription).
\end{cousinweb}

\begin{grammarlens}[title={What you've built}]
One more word for getting around.
\end{grammarlens}

\subsection*{Your turn}
\begin{itemize}
\raggedright
  \item \emph{Say it:} \emph{hōmen}
  \item \emph{Say it:} \emph{hōmen}, once more
  \item \emph{Say it:} \emph{minato hōmen no densha ni norimashita} — I took a train going toward the harbour.
  \item \emph{Recall:} say the Japanese for a suspension of service, then the Japanese for a timetable, then say \emph{hōmen} again
\end{itemize}

\begin{tcolorbox}[breakable,colback=teal!4,colframe=teal!35!black,title={Before you move on}]
What does \ja{ほうめん} mean? (\textquotedblleft{}A direction (of travel); an area\textquotedblright{}.) Say it once more.
\end{tcolorbox}
\section[\texorpdfstring{\ja{はなれる} (hanareru)}{はなれる (hanareru)}]{\ja{はなれる} (hanareru) — to move away; to be apart}

\label{lesson:JA-C377-hanareru}



\begin{quote}
Before the new one: say the Japanese for a timetable, then the Japanese for a direction (of travel).
\end{quote}

\subsection*{You'll want to know: \ja{はなれる}}
\textbf{\ja{はなれる}} — \emph{hanareru} — \textquotedblleft{}to move away; to be apart\textquotedblright{}.

\begin{quote}
\textbf{\ja{でんしゃが} \ja{えきを} \ja{はなれました。}} — \emph{densha ga eki o hanaremashita} — \textquotedblleft{}The train pulled away from the station.\textquotedblright{}
\end{quote}

Polite \textbf{\ja{はなれます}}, \emph{hanaremasu}; an ichidan verb. Past (polite): \textbf{\ja{ふねが} \ja{みなとを} \ja{はなれました。}}, \emph{fune ga minato o hanaremashita}.

\begin{cousinweb}
\textbf{\ja{はなれる}} (\emph{hanareru}) is a native Japanese word (wago), inherited rather than borrowed from Chinese; its kanji was attached later as a kun'yomi (JMdict).
\end{cousinweb}

\begin{grammarlens}[title={What you've built}]
One more word for getting around.
\end{grammarlens}

\subsection*{Your turn}
\begin{itemize}
\raggedright
  \item \emph{Say it:} \emph{hanareru}
  \item \emph{Say it:} \emph{hanareru}, once more
  \item \emph{Say it:} \emph{densha ga eki o hanaremashita} — The train pulled away from the station.
  \item \emph{Recall:} say the Japanese for a timetable, then the Japanese for a direction (of travel), then say \emph{hanareru} again
\end{itemize}

\begin{tcolorbox}[breakable,colback=teal!4,colframe=teal!35!black,title={Before you move on}]
What does \ja{はなれる} mean? (\textquotedblleft{}To move away; to be apart\textquotedblright{}.) Say it once more.
\end{tcolorbox}
\section[(dialogue)]{Looking back: from In Front of the Station to Delays and Cancellations}

\label{lesson:JA-R377-review-372-377}



\begin{quote}
Say the Japanese for in front of a station and a kiosk.
\end{quote}

\subsection*{Your turn}
Say these aloud:

\begin{itemize}
\raggedright
  \item in front of a station, a kiosk, to meet up, a waiting room, a vending machine
  \item ticket inspection, a commuter pass, a book of tickets, fare adjustment, to pass by
  \item a railway track, a level crossing, stopping (of a train or bus), departure (of a train or bus), to go near
  \item the first train of the day, the last train of the day, the terminus, to ride past one's stop, to get on board
  \item the bullet train, a rapid train, a private railway, a route, to skip
  \item a delay, a suspension of service, a timetable, a direction (of travel), to move away
\end{itemize}

\begin{tcolorbox}[breakable,colback=teal!4,colframe=teal!35!black,title={Before you move on}]
In front of a station? (\textbf{\ja{えきまえ}}, \emph{ekimae}.) A kiosk? (\textbf{\ja{ばいてん}}, \emph{baiten}.) To meet up? (\textbf{\ja{まちあわせる}}, \emph{machiawaseru}.) A waiting room? (\textbf{\ja{まちあいしつ}}, \emph{machiaishitsu}.) A vending machine? (\textbf{\ja{じどうはんばいき}}, \emph{jidōhanbaiki}.) Ticket inspection? (\textbf{\ja{かいさつ}}, \emph{kaisatsu}.) A commuter pass? (\textbf{\ja{ていきけん}}, \emph{teikiken}.) A book of tickets? (\textbf{\ja{かいすうけん}}, \emph{kaisūken}.) Fare adjustment? (\textbf{\ja{せいさん}}, \emph{seisan}.) To pass by? (\textbf{\ja{とおりすぎる}}, \emph{tōrisugiru}.) A railway track? (\textbf{\ja{せんろ}}, \emph{senro}.) A level crossing? (\textbf{\ja{ふみきり}}, \emph{fumikiri}.) Stopping (of a train or bus)? (\textbf{\ja{ていしゃ}}, \emph{teisha}.) Departure (of a train or bus)? (\textbf{\ja{はっしゃ}}, \emph{hassha}.) To go near? (\textbf{\ja{ちかよる}}, \emph{chikayoru}.) The first train of the day? (\textbf{\ja{しはつ}}, \emph{shihatsu}.) The last train of the day? (\textbf{\ja{しゅうでん}}, \emph{shūden}.) The terminus? (\textbf{\ja{しゅうてん}}, \emph{shūten}.) To ride past one's stop? (\textbf{\ja{のりこす}}, \emph{norikosu}.) To get on board? (\textbf{\ja{のりこむ}}, \emph{norikomu}.) The bullet train? (\textbf{\ja{しんかんせん}}, \emph{shinkansen}.) A rapid train? (\textbf{\ja{かいそく}}, \emph{kaisoku}.) A private railway? (\textbf{\ja{してつ}}, \emph{shitetsu}.) A route? (\textbf{\ja{ろせん}}, \emph{rosen}.) To skip? (\textbf{\ja{とばす}}, \emph{tobasu}.) A delay? (\textbf{\ja{ちえん}}, \emph{chien}.) A suspension of service? (\textbf{\ja{うんきゅう}}, \emph{unkyū}.) A timetable? (\textbf{\ja{じこくひょう}}, \emph{jikokuhyō}.) A direction (of travel)? (\textbf{\ja{ほうめん}}, \emph{hōmen}.) To move away? (\textbf{\ja{はなれる}}, \emph{hanareru}.)
\end{tcolorbox}
